\chapter{作品概述}
\section{功能与特性}
限 400 字以内，可加入重要图片。说明作品的主要功能、使用方式和关键特性。
\section{应用领域}
限 400 字以内，可加入重要图片。说明目标用户、应用场景及实际需求。
\section{主要技术特点}
限 400 字以内。说明 FPGA 在系统中的作用、核心算法、并行处理结构、接口及软硬件协同方式。
\section{主要性能指标}
限 200 字以内，建议用表格列出目标值、实测值及测试条件。
\begin{table}[htbp]
\centering
\caption{主要性能指标（填写示例）}
\label{tab:performance}
\begin{tabularx}{\textwidth}{XXXX}
\toprule
指标 & 目标值 & 实测值 & 测试条件 \\
\midrule
待填写 & 待填写 & 待填写 & 待填写 \\
\bottomrule
\end{tabularx}
\end{table}
\section{主要创新点}
限 200 字以内，逐点给出，说明创新对应的问题和验证依据。
\begin{enumerate}
\item 创新点一：待填写。
\item 创新点二：待填写。
\end{enumerate}
\section{设计流程}
限 200 字以内，可加入重要图片。概述需求分析、方案选择、模块设计、仿真验证、系统集成及测试迭代过程。
